\uuid{hgag}
\titre{Calcul de dérivée}
\auteur{bodin ; hanani ; mzari}
\organisation{exo7}
\licence{CC BY-NC-SA 4.0}
\datecreate{2020-01-20}
\niveau{L1}
\module{Analyse}
\chapitre{Dérivabilité des fonctions réelles}
\sousChapitre{Autre}
\difficulte{1}
\theme{Dérivées, calcul de dérivée}
\exo7qcmid{lille-1.derivees.009}

\contenu{
\begin{question}
\qclassification{124.00}
\qtags{motcle=calcul de dérivée}
Soit $\displaystyle f(x)=x^4-x^3+1$. Quelles sont les bonnes réponses ?
\begin{answers}
    \good{$f$ admet un minimum local au point $\displaystyle \frac{3}{4}$.}
    \bad{$f$ admet un maximum local au point $0$.}
    \bad{$f$ admet un minimum local au point $0$.}
    \good{$f$ admet un point d'inflexion au point $0$.}
\end{answers}
\begin{explanations}
On a $f'\left(\frac{3}{4}\right)=0$ et $f''\left(\frac{3}{4}\right)>0$. Donc $f$ admet un minimum au point $\displaystyle \frac{3}{4}$. On vérifie aussi que $f''$ s'annule en $0$ en changeant de signe. Donc $f$ admet un point d'inflexion au point $0$.
\end{explanations}
\end{question}
}
